\section{Scaling macro: de las curvas de los modelos a las señales de capital y de precios}

La expansión de la IA suele resumirse en una sola curva, pero el sistema industrial produce al menos cuatro tipos de señales a la vez: la señal técnica, que muestra cómo cambia la capacidad de los modelos según los recursos de entrenamiento; la señal de producto, que es el crecimiento en el uso de las herramientas; la señal de capital, que es la inversión en centros de datos y chips; y la señal de oferta y demanda, que dan los precios de renta de GPU. Estas señales pueden respaldarse entre sí, pero no son lo mismo. El progreso técnico no garantiza ingresos comerciales, el crecimiento de los ingresos no garantiza que el capital se asigne con eficacia y la inversión de capital tampoco garantiza que cada pieza de hardware se aproveche de manera eficiente.

\subsection{La previsibilidad de la capacidad de los modelos}

El valor de ingeniería de la scaling law está en la gestión del presupuesto y del riesgo: con experimentos anteriores, un equipo puede estimar cómo cambiará la pérdida si aumentan los datos, los parámetros o el cómputo, y con eso decidir el tamaño del siguiente entrenamiento. Mientras más lejos llega el intervalo de predicción, mayor es el error que introducen los cambios de arquitectura, la calidad de los datos, el posentrenamiento y la deriva de la evaluación. Las gráficas macro sirven para mostrar la dirección a largo plazo, pero no reemplazan la auditoría de cada modelo, cada tarea y cada recipe de entrenamiento.

\begin{figure}[H]
\centering
\includegraphics[width=0.94\textwidth]{images/00-44-00-macro-ai-scaling.jpg}
\caption{El ponente usa una gráfica macro de la expansión de los modelos para mostrar que el AI scaling presenta regularidades observables en varios periodos históricos.}
\figsource{Video oficial de Stanford Online 00:44:00, `images/00-44-00-macro-ai-scaling.jpg`.}
\end{figure}

\begin{knowledgebox}{Lectura de la figura: distinguir puntos observados, línea de ajuste y supuestos sobre el futuro}
Los puntos observados provienen de proyectos de entrenamiento ya publicados o estimados; la línea de ajuste resume la tendencia pasada, y el tramo que se extiende hacia la derecha incluye supuestos sobre el futuro. Conviene revisar primero las unidades de los ejes y la fuente de los datos, y después preguntar si las generaciones de modelos son comparables, si el posentrenamiento se tomó en cuenta y si la evaluación cambió. Un buen ajuste solo demuestra que los datos elegidos muestran una regularidad clara en ese sistema de coordenadas; no demuestra que la siguiente generación de sistemas vaya a crecer con la misma pendiente.
\end{knowledgebox}

\subsection{Adopción del producto: qué indica el crecimiento de Claude Code}

Las señales de producto están más cerca de la demanda real que los benchmark. El ponente muestra el crecimiento de los commits diarios de Claude Code a lo largo de doce meses para ilustrar que los coding agent se están incorporando rápidamente al proceso de desarrollo. Este indicador significa, como mínimo, que los usuarios están dispuestos a subir al repositorio código generado o asistido por el modelo, pero por sí solo no responde sobre la calidad del código, el costo del retrabajo, los incidentes de seguridad, la retención de clientes de pago ni la productividad neta de distintos equipos.

\begin{figure}[H]
\centering
\includegraphics[width=0.94\textwidth]{images/00-46-02-claude-code-growth.jpg}
\caption{Los commits diarios de Claude Code crecieron rápidamente en doce meses, lo que muestra la velocidad con la que los productos de agent entran en flujos de trabajo reales.}
\figsource{Video oficial de Stanford Online 00:46:02, `images/00-46-02-claude-code-growth.jpg`.}
\end{figure}

\begin{importantbox}{Lectura de la figura: el volumen de adopción no es una conclusión sobre productividad}
El número de commits puede demostrar que la intensidad de uso aumenta, pero no demuestra directamente que cada commit eleve la producción neta. Una evaluación completa también requiere merge rate, review time, tasa de defectos, tasa de reversiones, cobertura de desarrolladores y costo de mantenimiento a largo plazo. El crecimiento del producto sigue siendo importante porque amplía las fuentes de context y de retroalimentación; solo que esa retroalimentación debe pasar por filtros de calidad y de permisos antes de alimentar la siguiente ronda de mejoras del modelo o de las herramientas.
\end{importantbox}

\subsection{El gasto de capital de la Great Transition}

Cuando la demanda de aplicaciones, el tamaño de los modelos y las expectativas de capacidad de cómputo suben al mismo tiempo, los proveedores de nube y las empresas de infraestructura compran por adelantado terrenos, energía eléctrica, redes y aceleradores. \term{Capex} es capital expenditure, es decir, el gasto de capital destinado a activos de largo plazo; a diferencia de los gastos operativos cotidianos, suele tener ciclos de construcción largos, recuperación lenta y una salida difícil. Si los pronósticos de demanda de IA se desvían, tanto la sobreconstrucción como la escasez de oferta pueden prolongarse durante años.

\begin{figure}[H]
\centering
\includegraphics[width=0.94\textwidth]{images/00-46-45-great-transition.jpg}
\caption{El ponente describe el aumento reciente del gasto de capital en infraestructura como una señal macro de la Great Transition.}
\figsource{Video oficial de Stanford Online 00:46:45, `images/00-46-45-great-transition.jpg`.}
\end{figure}

\begin{knowledgebox}{Lectura de la figura: el capex es una apuesta, no un valor ya realizado}
La curva de gasto de capital indica que las empresas creen que habrá suficientes ingresos de software y valor estratégico para sostener la inversión en hardware. Una vez construido el centro de datos, su valor depende además de la tasa de utilización, la eficiencia de los modelos, el precio de la electricidad, la capacidad de interconexión y la demanda de los clientes. Es razonable tomar el capex como una expresión cuantitativa de la confianza en la demanda; tomarlo como prueba de que la productividad o el beneficio social ya se materializaron implica saltarse una cadena que todavía no se ha verificado.
\end{knowledgebox}

\subsection{Precios de renta de GPU: el termómetro de un mercado heterogéneo de cómputo}

El precio resulta atractivo porque condensa en un solo número la oferta, la demanda, la región, la energía eléctrica, la generación del hardware y el plazo del contrato. El ponente muestra la evolución de los precios de renta de GPU observada por AMP para ilustrar que la escasez no aumenta de forma monótona: la entrada de nueva oferta, la depreciación de las tarjetas antiguas, las mejoras en la eficiencia del software y los cambios en la estructura de la demanda modifican el precio. Más importante aún, las distintas GPU no son una misma mercancía; el precio de referencia debe leerse junto con la memoria de video, la interconexión, la confiabilidad, el tiempo de disponibilidad y el tamaño del clúster.

\begin{figure}[H]
\centering
\includegraphics[width=0.94\textwidth]{images/00-47-45-gpu-rental-prices.jpg}
\caption{Los precios de renta de GPU cambian según la generación del hardware y la etapa del mercado, y reflejan el efecto conjunto de la oferta y la demanda, la depreciación y la sustituibilidad.}
\figsource{Video oficial de Stanford Online 00:47:45, `images/00-47-45-gpu-rental-prices.jpg`.}
\end{figure}

\begin{warningbox}{Lectura de la figura: el precio por GPU-hour sigue sin ser el precio de una mercancía uniforme}
Aunque todo se exprese por hora, una sola tarjeta, un nodo de ocho tarjetas y un clúster grande sin bloqueo tienen valores completamente distintos; la renta de corto plazo, la reservada, la interrumpible y los contratos de largo plazo tampoco se pueden comparar directamente. El entrenamiento depende además de la topología de red, el rendimiento del almacenamiento, la pila de software y la recuperación ante fallas. La gráfica de precios puede mostrar una tendencia, pero sin las especificaciones y las condiciones del servicio no puede demostrar que «el cómputo ya es barato» ni que «cierta tarjeta conviene más con seguridad».
\end{warningbox}

\subsection{Resumen de la sección}

Las curvas de los modelos, la adopción de productos, el gasto de capital y los precios de renta responden, respectivamente, a preguntas de tecnología, demanda, inversión y oferta y demanda. Ponerlas lado a lado permite observar si la transición ocurre a través de varias capas, pero no permite fundir los cuatro tipos de señales en una cadena causal necesaria. La siguiente sección recurre a los ciclos históricos de infraestructura para buscar patrones reutilizables y, al mismo tiempo, deja claros los límites de la analogía.
